\section{Evaluation framework}
\label{sec:framework}

\paragraph{Models and references.}
Tab.~\ref{tab:models} lists what each model receives and produces.
Two references bound every score: the \emph{real} video of the same actions (ceiling) and the first frame held \emph{still} for 16 frames (floor: perfect image quality, no motion).
Cosmos videos are resampled to the IRASim format (16 frames, $256\times320$) before scoring; for the base model we keep every third of 49 frames at 16\,fps, which spans the same 3\,s.

\begin{table*}[t]
\centering
\footnotesize
\begin{tabular}{@{}llll@{}}
\toprule
Model & Conditioning & Native output & Steps \\
\midrule
IRASim-Frame-Ada & frame 0 + 15 actions, per frame & $256\times320$, 16 f & 50 \\
IRASim-Video-Ada & frame 0 + 15 actions, per video & $256\times320$, 16 f & 50 \\
LVDM & frame 0 + actions, per video & $256\times320$, 16 f & 50 \\
VDM (pixels) & frame 0 + actions, per video & $256\times320$, 16 f & 50 \\
Cosmos-P2.5 action & frame 0 + chunks of 12 actions & $256\times320$, 13 f per chunk & 35 \\
Cosmos-P2.5 base & frame 0 + instruction text & $480\times640$, 49 f, 16 fps & 35 \\
\bottomrule
\end{tabular}
\caption{Models compared on Bridge. Cosmos action: two chunks, frames 0 to 15 kept; guidance 0. Cosmos base: guidance 7, 24 core clips only (3 min per clip on an A100).}
\label{tab:models}
\end{table*}

\paragraph{Metric groups.}
Tab.~\ref{tab:metrics} states for each group what it measures and what it cannot establish.
All are computed per clip and compared with paired Wilcoxon tests over the same clips.

\begin{table*}[t]
\centering
\footnotesize
\begin{tabular}{@{}p{3.1cm}p{5.4cm}p{6.4cm}@{}}
\toprule
Group (metrics) & Measures & Cannot establish \\
\midrule
Visual quality (PSNR, SSIM, LPIPS vs the real clip) & how close each frame looks to the recorded frame & whether the robot moved as commanded: a still video keeps high SSIM; small spatial shifts are punished like large errors \\
Temporal consistency (DINOv2 \cite{oquab2023dinov2} similarity of adjacent frames and to frame 0; optical-flow warp error) & smoothness and identity stability across frames & correct motion: the still video scores perfectly \\
2-D trajectory (CoTracker3 \cite{karaev2024cotracker3} end-effector track vs the real clip; NDTW \cite{ilharco2019ndtw}, error in pixels, path ratio) & whether the gripper moves along the same image path as in the real clip & depth and orientation; needs the real clip; relies on picking the end-effector point (VLM-checked) \\
Physical plausibility and outcome (two VLM judges) & violations seen by a judge (objects vanish or change, robot deforms) and whether the final state matches the real clip & reliability: judges differ in strictness and change answers between runs \\
3-D action consistency (Cosmos3-Nano inverse dynamics, ours) & whether the 3-D end-effector motion read from the video matches the commanded actions, in cm; needs no real video & object motion and contact; bounded below by the inverse dynamics model's own error (1.50\,cm on real clips) \\
\bottomrule
\end{tabular}
\caption{Metric groups. Task completion is covered by the outcome judge and, for executability, by replay in simulation (Sec.~\ref{sec:improve}).}
\label{tab:metrics}
\end{table*}
